\section{The local limit}\label{sec:local}

In this section $0<\alpha\le f\le\beta<\infty$, $c=c_f$ and $\lambda=m/c$. By Theorem~\ref{thm:normal},
$S_t/t\to c$ almost surely; since $\alpha\le S_t/t\le\beta$, for every $\eps>0$
\begin{equation}\label{eq:tailsup}
   \E\sup_{t\ge\eps n}\bigl|S_t/t-c\bigr|\to0\qquad(n\to\infty).
\end{equation}

\subsection{Ball shapes and the limiting probabilities}

Fix $r\ge1$. A \emph{ball shape} is a finite rooted directed simple tree $H$ with vertex set $V$ and root $o$ in
which every vertex is within distance $r$ of $o$ and every vertex at distance $r$ is a leaf. Let $I$ be the set
of vertices at distance less than $r$ (the interior), and for $j\in V$ let $q_j$ and $k_j$ be the out-degree and
the in-degree of $j$ in $H$. We write $j\to i$ for an edge sent by $j$ and received by $i$. The shape is
\emph{admissible} if $q_j=m$ for every $j\in I$ and $q_j\le m$ for every $j$. Let $A(H)$ be the number of
automorphisms of $H$ that fix the root and preserve orientations, and $(m)_q=m!/(m-q)!$.

An \emph{age vector} is $x\in(0,1)^V$ with distinct coordinates; it is \emph{compatible} if $x_j>x_i$ for every
edge $j\to i$. Given a compatible age vector, let $(Y_i)_{i\in V}$ be independent processes, where $Y_i$ is
defined on $[x_i,1]$, starts at $0$, jumps by $+1$ at rate $\lambda f(Y_i(t))/t$ at time $t$, and in addition
jumps by $+1$ at each of the times $x_j$, $j\to i$. Define
\begin{equation}\label{eq:Phi}
  \Phi_H(x)=\prod_{j\in V}(m)_{q_j}\ \prod_{j\to i}\frac{1}{c\,x_j}\ \,
     \E\Bigl[\,\prod_{j\to i}f\bigl(Y_i(x_j-)\bigr)\prod_{i\in I}\mathbf 1_{\{Y_i(1)=k_i\}}\Bigr]
\end{equation}
for compatible $x$ and $\Phi_H(x)=0$ otherwise, and
\begin{equation}\label{eq:PH}
   \mathbf P_H=\frac{1}{A(H)}\int_{(0,1)^V}\Phi_H(x)\,\mathrm dx .
\end{equation}

\begin{theorem}\label{thm:localbdd}
For every $r\ge1$ and every admissible ball shape $H$ of radius $r$,
\[
   \frac1n\,\#\bigl\{u\in[n]:\ B_r(G_n,u)\cong H\bigr\}\ \longrightarrow\ \mathbf P_H\qquad\text{in probability,}
\]
and the proportion of $u$ whose ball $B_r(G_n,u)$ is not an admissible ball shape tends to $0$ in probability.
Consequently $G_n$ converges locally in probability to the random rooted tree whose balls have the laws
$(\mathbf P_H)$; this tree is infinite, simple, and all its vertices have out-degree $m$.
\end{theorem}

The expectation in \eqref{eq:Phi} is a product of explicit factors.

\begin{proposition}\label{prop:factors}
Let $X$ be the pure birth process started at $0$ with rates $\lambda f(k)$, and $h_\lambda(u)=\E f(X(u))$. For
compatible $x$ the expectation in \eqref{eq:Phi} equals $\prod_{i\in V}F_i$, where:
for $i\in I$ with children of ages $t_1<\dots<t_k$ $(k=k_i)$, and $t_0=x_i$, $t_{k+1}=1$,
\[
   F_i=\prod_{\ell=0}^{k-1}f(\ell)\ \prod_{\ell=0}^{k}\Bigl(\frac{t_\ell}{t_{\ell+1}}\Bigr)^{\lambda f(\ell)};
\]
for $i\notin I$, $F_i=1$ if $k_i=0$ and $F_i=h_\lambda(\log(t_1/x_i))$ if $k_i=1$, where $t_1$ is the age of the
child of $i$.
\end{proposition}

\begin{proof}
The processes are independent, so the expectation factorises. For $i\in I$ the event $Y_i(1)=k_i$ says that
$Y_i$ makes no jump other than the $k_i$ forced ones; on this event $Y_i=\ell$ on $[t_\ell,t_{\ell+1})$, the
probability of no other jump on that interval is $\exp(-\lambda f(\ell)\int_{t_\ell}^{t_{\ell+1}}\mathrm dt/t)$,
and the weight at the $\ell$th forced jump is $f(\ell-1)$. A vertex $i\notin I$ is a leaf of $H$ with no
constraint at time $1$; if its neighbour is its child of age $t_1$, then $Y_i(t_1-)$ has the law of
$X(\log(t_1/x_i))$.
\end{proof}

For instance, if $H$ is the star in which the root has no child, the root contributes $x_o^{\lambda f(0)}$, and
each of its $m$ parents, of age $y<x_o$, contributes $h_\lambda(\log(x_o/y))/(c\,x_o)$, which integrates to
$\int_0^\infty e^{-u}h_\lambda(u)\,\mathrm du/c=1$ over $y\in(0,x_o)$ by Lemma~\ref{lem:scalar}(ii); since
$(m)_m=A(H)=m!$, this gives $\mathbf P_H=\int_0^1x^{\lambda f(0)}\,\mathrm dx=p_0(c)$, in agreement with
Theorem~\ref{thm:normal}. The same computation shows that the age $y$ of a parent of a vertex of age $x$ has
density $h_\lambda(\log(x/y))/(cx)$ on $(0,x)$, which is not uniform unless $f$ is constant.

\subsection{Forced slots}

Fix an admissible $H$ with $e=|V|-1$ edges, an injective labelling $\iota:V\to[n]$ with $\iota(j)>\iota(i)$ for
every edge $j\to i$, and a \emph{slot assignment} $\sigma$, which gives to every edge $j\to i$ a slot of
$\iota(j)$, different edges receiving different slots. There are $\prod_j(m)_{q_j}$ slot assignments. Let
$E_{\iota,\sigma}$ be the event that for every edge $j\to i$ the assigned slot of $\iota(j)$ chooses
$\iota(i)$, and let $D_\iota=\{D_{\iota(i)}(n)=k_i\ \text{for all }i\in I\}$.

The \emph{forced process} is the graph process in which every assigned slot chooses its prescribed endpoint with
probability one and all other slots follow the original rule. Write $\Prob^{\sigma}$, $\E^\sigma$ for its law.
Multiplying the transition probabilities along a history gives, for every nonnegative function $\Psi$ of the
history up to time $n$,
\begin{equation}\label{eq:com}
   \E\bigl[\Psi\,\mathbf 1_{E_{\iota,\sigma}}\bigr]
   =\E^\sigma\Bigl[\Psi\prod_{j\to i}\frac{f(D^-_{j\to i})}{S^-_{j\to i}}\Bigr],
\end{equation}
where $D^-_{j\to i}$ is the in-degree of $\iota(i)$ and $S^-_{j\to i}$ the total weight of the eligible vertices
immediately before the slot assigned to $j\to i$.

\begin{lemma}[The forced process has the same normalisation]\label{lem:forcedS}
Suppose that all labels $\iota(j)$ are at least $\eps n$. Under $\Prob^\sigma$,
$\sup_{\eps n\le t\le n}|S_t/t-c|\to0$ in $L^1$ as $n\to\infty$, uniformly in $\iota$ and $\sigma$.
\end{lemma}

\begin{proof}
Couple the forced process with the original one: they coincide until the first assigned slot, and afterwards,
at each slot that is not assigned, the two endpoints are chosen by a maximal coupling of the two conditional
laws. Let $\bar K$ be the number of labels at which the two in-degree sequences have differed at some time. If
the two weight vectors at a slot of vertex $t+1$ are $w$ and $w'$ with totals $W,W'\ge\alpha t$, then
\[
   \bigl\|w/W-w'/W'\bigr\|_1\le\frac{2\|w-w'\|_1}{W}\le\frac{2\beta\bar K}{\alpha t},
\]
so the two choices differ with probability at most $C\bar K/t$, in which case $\bar K$ increases by at most $2$;
at an assigned slot $\bar K$ increases by at most $2$. Hence
$\E\bar K\le 2e\exp(2mC\sum_{t\ge\eps n/2}^{n}1/t)\le K_\eps$, a constant that does not depend on $n$. Since the
two values of $S_t$ differ by at most $\beta\bar K$, the claim follows from \eqref{eq:tailsup}.
\end{proof}

\begin{lemma}[Tagged degrees]\label{lem:tagged}
Let $\iota=\iota_n$ and $\sigma$ be as above, and suppose that $\iota_n(j)/n\to x_j$ for all $j$, where $x$ is a
compatible age vector with all $x_j\ge\eps$. Under $\Prob^\sigma$, the vector formed by the in-degrees
$D^-_{j\to i}$ $(j\to i)$ and $D_{\iota(i)}(n)$ $(i\in V)$ converges in total variation to the vector formed by
$Y_i(x_j-)$ and $Y_i(1)$.
\end{lemma}

\begin{proof}
The vertices $\iota(i)$, $i\in V$, are called tagged. Consider a slot of a vertex $t+1$ that is not assigned.
Given the history, the tagged vertex $\iota(i)\le t$ is chosen with probability $f(D_{\iota(i)})/S$, where $S$
is the current total weight, and at most one tagged vertex is chosen. Compare with independent Bernoulli
variables of parameters $f(D_{\iota(i)})/(ct)$: the total variation distance between the two laws of the
increment of the tagged vector is at most
\[
   K_1\,\frac{|S/t-c|}{t}+\frac{K_1}{t^2},
\]
with $K_1$ depending on $|V|,m,\alpha,\beta$ only, because $S\ge\alpha t$, $c\ge\alpha$ and the probability
that two of the Bernoulli variables equal one is $O(t^{-2})$. At an assigned slot the prescribed tagged vertex
gains one and the others gain nothing, in both systems. Using maximal couplings slot by slot, the tagged vector
of the forced process coincides up to time $n$ with a system of independent chains, one for each tagged vertex,
which start at $0$ at the birth of the vertex, jump at each unassigned slot of a vertex $t+1$ with probability
$f(\text{state})/(ct)$, and jump at their assigned slots, except on an event of probability at most
\[
   K_1m\log(1/\eps)\ \E^\sigma\!\sup_{t\ge\eps n}|S_t/t-c|+O\bigl(1/(\eps n)\bigr),
\]
which tends to $0$ by Lemma~\ref{lem:forcedS} (within a batch $S$ differs from $S_t$ by at most $2m\beta$).

It remains to compare each independent chain with $Y_i$. Over the slots of vertex $t+1$ the chain in state $k$
jumps exactly once with probability $mf(k)/(ct)+O(t^{-2})$ and more than once with probability $O(t^{-2})$;
the process $Y_i$ in state $k$ has the same two estimates on the age interval $[t/n,(t+1)/n]$, since
$\int\lambda f(k)\,\mathrm ds/s=\lambda f(k)\log(1+1/t)$. The total variation distance between the two paths,
observed at the ages $t/n$, is therefore $O(\sum_{t\ge\eps n}t^{-2})=O(1/(\eps n))$, apart from the effect of the
forced times, which are $\iota_n(j)/n$ for the chain and $x_j$ for $Y_i$: moving a forced time by $\delta$
changes the law of the observed vector by at most the probability of an unforced jump in a window of length
$\delta$, which is $O(\delta/\eps)$. Finally the in-degree of a target immediately before its assigned slot
differs from its in-degree at the beginning of the batch with probability $O(1/n)$.
\end{proof}

\begin{proposition}\label{prop:embed}
Under the assumptions of Lemma~\ref{lem:tagged},
\[
   n^{e}\,\Prob\bigl(E_{\iota_n,\sigma}\cap D_{\iota_n}\bigr)\ \longrightarrow\
   \prod_{j\to i}\frac1{c\,x_j}\ \E\Bigl[\prod_{j\to i}f\bigl(Y_i(x_j-)\bigr)\prod_{i\in I}\mathbf 1_{\{Y_i(1)=k_i\}}\Bigr],
\]
and the left-hand side is bounded by $(4C/\eps)^e$ for all labellings with labels at least $\eps n$ and $n$
large. The convergence is uniform on every compact set of compatible age vectors.
\end{proposition}

\begin{proof}
Apply \eqref{eq:com} with $\Psi=\mathbf 1_{D_\iota}$. Each factor $nf(D^-)/S^-$ is at most
$n\beta/(\alpha(\eps n-1))$, and $S^-_{j\to i}/n\to c\,x_j$ in $L^1$ under $\Prob^\sigma$ by
Lemma~\ref{lem:forcedS}. The numerators and the indicator are functions of the vector of
Lemma~\ref{lem:tagged}. The limit is continuous in $x$ on the set of compatible age vectors, and convergence
along every sequence $\iota_n(\cdot)/n\to x$ gives the uniformity.
\end{proof}

\subsection{Proof of Theorem~\ref{thm:localbdd}}

\emph{Counting.} Let $\mathcal A_\iota$ be the event that $\iota$ is an isomorphism from $H$ onto
$B_r(G_n,\iota(o))$. A ball isomorphic to $H$ has exactly $A(H)$ such isomorphisms, so
\begin{equation}\label{eq:count}
   \Prob\bigl(B_r(G_n,U_n)\cong H\bigr)=\frac{1}{n\,A(H)}\sum_\iota\Prob(\mathcal A_\iota).
\end{equation}
We compare $\mathbf 1_{\mathcal A_\iota}$ with $\sum_\sigma\mathbf 1_{E_{\iota,\sigma}\cap D_\iota}$. Every edge
of $H$ has an endpoint in $I$. On $D_\iota$ no edge $j\to i$ with $i\in I$ can be doubled, because the in-degree
of $\iota(i)$ is exactly the number of its children; and an edge $j\to i$ with $j\in I$ cannot be doubled on
$E_{\iota,\sigma}$, because all $m$ slots of $\iota(j)$ are assigned to distinct endpoints. Hence at most one
$\sigma$ contributes, and on $\mathcal A_\iota$ exactly one does. Conversely, on $E_{\iota,\sigma}\cap D_\iota$
all edges of $G_n$ at the vertices $\iota(i)$, $i\in I$, are images of edges of $H$; by induction on the
distance, the vertices at distance $d\le r$ from $\iota(o)$ in $G_n$ are the images of the vertices at distance
$d$ in $H$. Thus $\mathcal A_\iota$ can fail only if $G_n$ has an edge between the images of two vertices at
distance $r$. If all labels are at least $\eps n$, Lemma~\ref{lem:edges} bounds the probability of
$E_{\iota,\sigma}$ together with one such extra edge by $(4mC/(\eps n))^{e+1}$. There are at most
$n^{|V|}=n^{e+1}$ labellings, so
\begin{equation}\label{eq:count2}
   \Prob\bigl(B_r(G_n,U_n)\cong H,\ B_r(G_n,U_n)\subset[\eps n,n]\bigr)
   =\frac{1}{A(H)}\ \frac{1}{n^{|V|}}\sum_{\iota\ge\eps n}\ \sum_\sigma n^{e}\,\Prob(E_{\iota,\sigma}\cap D_\iota)+O_\eps(1/n).
\end{equation}

\emph{Riemann sums.} By Proposition~\ref{prop:embed} the summands are bounded, and they converge uniformly on
compact sets of compatible age vectors; the limit does not depend on $\sigma$. Labellings for which two ages are
within $\delta$ form a proportion $O(\delta+1/n)$ of all labellings. Letting $n\to\infty$ and then $\delta\to0$,
the right-hand side of \eqref{eq:count2} converges to $A(H)^{-1}\int_{[\eps,1]^V}\Phi_H(x)\,\mathrm dx$.

\emph{The cutoff.} By Proposition~\ref{prop:balls} the left-hand side of \eqref{eq:count2} differs from
$\Prob(B_r(G_n,U_n)\cong H)$ by at most $(\sum_{l\le r}B^l)\sqrt\eps$. The integrals over $[\eps,1]^V$ increase as
$\eps$ decreases and are bounded by one, so $\Phi_H$ is integrable and
$\Prob(B_r(G_n,U_n)\cong H)\to\mathbf P_H$. Balls that are not admissible ball shapes have vanishing probability
by Proposition~\ref{prop:cycles} and Corollary~\ref{cor:tight}, and tightness shows that $(\mathbf P_H)_H$ is a
probability distribution for each $r$.

\emph{Two roots.} Let $U_n,U_n'$ be independent uniform vertices. Their balls of radius $r$ intersect with
probability at most $\E|B_{2r}(G_n,U_n)|/n=O(1/n)$. For two admissible shapes $H_1,H_2$ repeat the argument with
the disjoint union $H_1\sqcup H_2$, which has $|V_1|+|V_2|-2$ edges, and two root factors $1/n^2$. An edge
between the image of $H_1$ and the image of $H_2$ can only join vertices at distance $r$ from the two roots and
does not change either ball, so the counting step is unchanged. No forced jump links the two components, the
processes $Y_i$ are independent, and the limiting density is $\Phi_{H_1}\Phi_{H_2}$. Therefore
\[
   \Prob\bigl(B_r(G_n,U_n)\cong H_1,\ B_r(G_n,U_n')\cong H_2\bigr)\to\mathbf P_{H_1}\mathbf P_{H_2}.
\]
With $H_1=H_2=H$ this says that the second moment of $n^{-1}\#\{u:B_r(G_n,u)\cong H\}$ converges to
$\mathbf P_H^2$, and the first moment converges to $\mathbf P_H$; hence the convergence in probability. Since
balls with more than $M$ vertices have expected proportion at most $(\sum_{l\le r}B^l)/M$, the convergence for
each shape gives convergence in probability of the empirical average of every bounded function of the ball of
radius $r$, which is local convergence in probability. The properties of the limit were proved in
Corollary~\ref{cor:tight}. \qed
